\section{Microscopic lower denominators and exact-window minorization}
\label{sec:v51-roof-minorization}
The lower local law has a consequence not supplied by local total
variation alone: a positive arithmetic reference is dominated by the
original density, uniformly in the roof outside its null set. This
gives lower denominators for arbitrarily small measurable roof sets.
For fixed windows it yields a positive common part of the two
conditional roof laws, and hence a reverse likelihood bound.

Use $P_{m,R}^{n,k}$, $G_{m,R}^{n,k}$ and $\delta_m$ from
Section~\ref{sec:v51-positive-remainder}. Let $d>0$ be fixed and set
$\mathcal T_{m,R}^{n,k}(d)=\{u:G_{m,R}^{n,k}(u)\ge d\}$.
Neither this definition nor any result below assigns a positive
reference to a zero arithmetic residue.

\begin{corollary}[Lower denominators at the original discrete labels]
\label{cor:v51-lower-denominators}
For all sufficiently large $m$, uniformly in $R,n,k$,
\begin{equation}\label{eq:v51-roof-lower-denominator}
 p_{n,R}(k,m,u)\ge \frac{d}{2m^2}
       \quad\hbox{for almost every }u\in\mathcal T_{m,R}^{n,k}(d).
\end{equation}
For every Borel set $E\subset\mathcal T_{m,R}^{n,k}(d)$ of finite
Lebesgue measure,
\begin{align}
 &\nu_R^*\{K_{n,R}=k,N_{n,R}=m,T_{n,R}\in E\}\notag\\
 &\qquad\ge m^{-2}(1-\delta_m/d)
                  \int_E G_{m,R}^{n,k}(u)\,du
       \ge \frac{d}{2m^2}\operatorname{Leb}(E).
                                    \label{eq:v51-measurable-denominator}
\end{align}
The set $E$ may be chosen after $m$ and may have arbitrarily small
positive measure. No asymptotic upper bound on that event is asserted.
\end{corollary}
\begin{proof}
By \eqref{eq:v51-lower-local-law}, $P\ge G-\delta_m$ almost
everywhere. On $\{G\ge d\}$ this is at least
$(1-\delta_m/d)G$. Take $m$ sufficiently large that $\delta_m<d/2$
and integrate on $E$. The measure is the unchanged original source,
so the displayed probability is exactly the integral of $p$.
\end{proof}

Fix also $h>0$ and consider all target windows $I_t=[t,t+h]$ on
which $G_{m,R}^{n,k}\ge d$ almost everywhere. Define
\[
 e_m(h)=\sup_{R,n,k,t}h^{-1}\int_{I_t}|P_{m,R}^{n,k}-G_{m,R}^{n,k}|,
 \quad F=\int_{I_t}P_{m,R}^{n,k},\quad H=\int_{I_t}G_{m,R}^{n,k}.
\]
Thus $e_m(h)\to0$ by Theorem~\ref{thm:v49-raw-local-TV} and
$H\ge dh$. Let $\mathbb P_{m,R}^{n,k,t}$ be the law of the roof
conditioned on the unchanged exact event
\[
 K_{n,R}=k,\qquad N_{n,R}=m,\qquad T_{n,R}\in I_t,
\]
and let $\mathbb Q_{m,R}^{n,k,t}$ have density
$\one_{I_t}G_{m,R}^{n,k}/H$. The lower bound on the reference is
pointwise here; a bound only on its total mass would not suffice for
the domination asserted next.

\begin{theorem}[An asymptotically full common conditional measure]
\label{thm:v51-conditional-minorization}
For all sufficiently large $m$, set
\begin{equation}\label{eq:v51-minorization-factor}
 \alpha_m(d,h)=\frac{1-\delta_m/d}{1+e_m(h)/d}\in(0,1].
\end{equation}
Uniformly on the indicated target family,
\begin{equation}\label{eq:v51-conditional-minorization}
 \mathbb P_{m,R}^{n,k,t}\ge
                   \alpha_m(d,h)\mathbb Q_{m,R}^{n,k,t}
                       \quad\hbox{as measures},\qquad
 \alpha_m(d,h)\longrightarrow1.
\end{equation}
If $\alpha_m<1$ there is a probability $\mathbb S_{m,R}^{n,k,t}$
on the same roof interval such that exactly
\begin{equation}\label{eq:v51-conditional-mixture}
 \mathbb P=\alpha_m\mathbb Q+(1-\alpha_m)\mathbb S.
\end{equation}
If $\alpha_m=1$, the two conditional laws are equal.
In particular, defining reverse likelihood divergence by
\[
 D_\infty(Q\Vert P)=\log\operatorname*{ess\,sup}_{Q}\frac{dQ}{dP},
 \qquad
 D(Q\Vert P)=\int\log\frac{dQ}{dP}\,dQ,
\]
one has
\begin{equation}\label{eq:v51-reverse-likelihood}
 0\le D(\mathbb Q\Vert\mathbb P)
 \le D_\infty(\mathbb Q\Vert\mathbb P)
 \le-\log\alpha_m(d,h)\longrightarrow0.
\end{equation}
No forward likelihood bound or single-roof path bridge is included.
\end{theorem}
\begin{proof}
The lower local law gives $P\ge(1-\delta_m/d)G$ on $I_t$.
Also $F\le H+he_m(h)\le H(1+e_m(h)/d)$. Divide the first
inequality by this upper bound for $F$, which is positive by
Corollary~\ref{cor:v51-lower-denominators}. This proves
\eqref{eq:v51-conditional-minorization}. The difference measure
$\mathbb P-\alpha_m\mathbb Q$ is nonnegative and has mass
$1-\alpha_m$; normalize it to obtain
\eqref{eq:v51-conditional-mixture}. This residual is a decomposition
of the same conditional law, not a replacement conditioning event.

Domination implies $\mathbb Q\ll\mathbb P$ and
$d\mathbb Q/d\mathbb P\le\alpha_m^{-1}$, proving the
$D_\infty$ estimate. The relative entropy is nonnegative by convexity
and is at most the essential logarithmic bound. It is well defined:
the negative part of $r\log r$, with $r=dQ/dP$, is integrable
against the probability $P$, and $r\le\alpha_m^{-1}$ bounds the
positive part. Finally \eqref{eq:v51-minorization-factor} tends
to one, since $\delta_m$ and $e_m(h)$ tend to zero uniformly.
\end{proof}

\begin{remark}[A fixed-radius arithmetic specialization]
\label{rem:v51-arithmetic-positivity}
At fixed $R$, on a fixed central compact set and on labels with
$\mathfrak a_R(k,n,m)\ge a_0>0$, the Gaussian has a positive
minimum and the fixed-radius kernel approximation is uniform.
It therefore supplies a $d>0$ for all sufficiently large counts.
For a parameter-uniform family the stated lower condition on
$\mathcal L_{m,R}$ must be retained through arithmetic transitions.
No claim that $\mathfrak a_R$ is identically one, or that all labels
satisfy this lower condition, is needed.
\end{remark}

\paragraph{Density versions and the remaining bridge problem.}
The essential bound \eqref{eq:v51-roof-lower-denominator} is a
statement about a density, not a positive probability for the event
$T_{n,R}=u$. That event has probability zero. At Lebesgue points it
can be read using the differentiation representative of the original
density. It supplies neither a value at an arbitrarily designated
exceptional roof nor an asymptotic numerator for a path selector
conditioned at that roof. The two-sided density estimate and a
same-roof path numerator are still required for the original
pointwise bridge problem.

\paragraph{Why the two directions of likelihood are different.}
For example, on $[0,1]$ take the reference density one. Let
$H_j=2^{j^2}$, $w_j=(jH_j)^{-1}$ and
\[
 p_j(u)=\frac{1+H_j\one_{[0,w_j]}(u)}{1+1/j}.
\]
These are probability densities, their $L^1$ errors tend to zero,
and they dominate the reference by $(1+1/j)^{-1}$. Hence their
reverse $D_\infty$ tends to zero, while their essential heights
diverge. The contribution of the spike to forward relative entropy
is asymptotic to $j\log2$, so forward relative entropy diverges.
This example is a check on the logical scope of the conclusions,
not a model for a billiard singularity. It explains why neither the
minorization nor the local variation theorem closes
\eqref{eq:v51-positive-criterion}.
